% Guion F02 ampliado a6:45 por solicitud del usuario el19/11/2026.
% Esta revisión edita el documento LaTeX abierto; conserva responsables y diapositivas.
% Codificación UTF-8. Compilar con pdfLaTeX (por ejemplo, en Overleaf).
\documentclass[11pt,a4paper]{article}
\usepackage[utf8]{inputenc}
\usepackage[T1]{fontenc}
\usepackage[spanish,es-nodecimaldot]{babel}
\usepackage[a4paper,margin=2.2cm,headheight=15pt]{geometry}
\usepackage{lmodern}
\usepackage{amsmath,amssymb}
\usepackage{xcolor}
\usepackage{parskip}
\usepackage{fancyhdr}
\usepackage[hidelinks]{hyperref}
\definecolor{petroleo}{HTML}{166C65}
\definecolor{grisindicacion}{HTML}{52616B}
\pagestyle{fancy}
\fancyhf{}
\fancyhead[L]{\small Guion oral F02 -- Grupo 6}
\fancyhead[R]{\small Duración objetivo: 6:45}
\fancyfoot[C]{\thepage}
\setlength{\parindent}{0pt}
\setlength{\parskip}{0.7em}
\newenvironment{indicacion}{\begin{quote}\small\color{grisindicacion}}{\end{quote}}
\newcommand{\lamina}[4]{%
  \section*{\color{petroleo}Diapositiva #1: #2}%
  \textbf{Responsable:} #3\hfill\textbf{Intervalo:} #4\par
}
\hypersetup{pdftitle={Guion oral F02 por diapositiva},pdfauthor={Benjamin Garcia, Luis Noriega y Matias Stierling}}
\begin{document}

\lamina{1}{Decisión y objetivo}{Benjamín}{0:00--0:40}

\begin{indicacion}\textbf{Uso:} el texto normal se dice en voz alta. Las indicaciones en gris y las marcas de clic no se leen. Se conserva el reparto por diapositiva. Objetivo: 6:45 a ritmo conversacional, aproximadamente 130--140 palabras por minuto durante el discurso, con pausas breves para clics y relevos. Los intervalos son referencias para ensayar, no una duración medida.\end{indicacion}

\begin{indicacion}\emph{Al entrar en pantalla:} Fundación en 3D, título del grupo y autores. Nada más.\end{indicacion}


Buenos días. Somos el grupo 6 y buscamos colocar la fundación de El Rosal en el menor tiempo posible, respetando los límites térmicos.



\begin{indicacion}\textbf{\ensuremath{\blacktriangleright} CLIC 1} --- aparece: 10,61 h y «60 m\textsuperscript{3}/h \ensuremath{\cdot} 12 \ensuremath{^\circ}C \ensuremath{\cdot} 9 circuitos». \emph{Cuándo:} apenas termines el saludo.\end{indicacion}


Proponemos una operación continua a 60 metros cúbicos por hora, con hormigón a 12 grados y nueve circuitos de refrigeración. La colocación completa dura 10,61 horas.



\begin{indicacion}\textbf{\ensuremath{\blacktriangleright} CLIC 2} --- aparece: «0 / 24 excedencias» y los límites. \emph{Cuándo:} después de decir «10,61 horas».\end{indicacion}


Evaluamos 24 escenarios con viento y radiación: ninguno excede los límites en los 14 días exigidos ni hasta los 28 simulados. Ahora explicaremos qué hace posible ese resultado y bajo qué supuestos.



\clearpage

\lamina{2}{Por qué aparece el diferencial}{Benjamín}{0:40--1:35}

\begin{indicacion}\emph{Al entrar en pantalla:} Corte de la fundación en gris y, a la derecha, la ecuación de balance.\end{indicacion}


El problema físico comienza con la hidratación del cemento.



\begin{indicacion}\textbf{\ensuremath{\blacktriangleright} CLIC 1} --- aparece: mancha naranja «hidratación Q(a)». \emph{Cuándo:} después de presentar la hidratación.\end{indicacion}


Esa reacción libera calor dentro del hormigón. Como la fundación es masiva, el calor del interior tarda en alcanzar el exterior por conducción. Por eso la temperatura no es uniforme.



\begin{indicacion}\textbf{\ensuremath{\blacktriangleright} CLIC 2} --- aparece: punto caliente, punto S25 y la ecuación de \ensuremath{\Delta}T. \emph{Cuándo:} antes de «el interior conserva calor».\end{indicacion}


El interior conserva calor mientras el borde puede enfriarse. Medimos la diferencia entre la máxima del hormigón colocado y la mínima de S25, siempre en el mismo instante. S25 corresponde a puntos del hormigón a unos 25 milímetros bajo las superficies exteriores; no es la temperatura del suelo.



\begin{indicacion}\textbf{\ensuremath{\blacktriangleright} CLIC 3} --- aparece: flechas de borde, viento, sol, mantas y suelo. \emph{Cuándo:} al decir «en los bordes».\end{indicacion}


En los bordes actúan la convección por viento, el sol y la radiación hacia el cielo. Las mantas frenan ese intercambio. Representamos el suelo a 15 grados, sin resolver su calentamiento.



\clearpage

\lamina{3}{Cómo colocamos y refrigeramos}{Benjamín}{1:35--2:38}

\begin{indicacion}\emph{Al entrar en pantalla:} Solo el suelo y el título; el corte se construye con el primer clic.\end{indicacion}

\begin{indicacion}\textbf{\ensuremath{\blacktriangleright} CLIC 1} --- aparece: las 11 capas suben una por una y aparecen los dos frentes. \emph{Cuándo:} al empezar a hablar.\end{indicacion}


Colocamos once capas de hasta 40 centímetros con dos frentes de 30 metros cúbicos por hora. Es continuo: no esperamos que cada capa endurezca. La unión exige mantener la trabajabilidad; la ventana de 150 minutos es una condición por verificar.



\begin{indicacion}\textbf{\ensuremath{\blacktriangleright} CLIC 2} --- aparece: los 4 niveles de tubos, el agua circulando y las cifras de la derecha. \emph{Cuándo:} cuando digas «dentro van nueve circuitos».\end{indicacion}


Dentro van nueve circuitos en cuatro alturas, con 18 litros por minuto por ramal abierto. El agua se idealiza uniforme a 18 grados. El modelo extrae calor según la diferencia de temperatura, con un tope global de 15 kilowatts; no significa extraer siempre esa potencia.



\begin{indicacion}\textbf{\ensuremath{\blacktriangleright} CLIC 3} --- aparece: gráfico de la regla 19/18 \ensuremath{^\circ}C y protección 504 h + sombra 90 \%. \emph{Cuándo:} antes de explicar la regla de apertura y las protecciones.\end{indicacion}


Cada circuito abre cuando está cubierto y su mínima monitoreada alcanza 19 grados; cierra a 18. Mantenemos la última manta hasta 504 horas desde la colocación de cada capa. Después, sombra del 90 por ciento y control térmico continuo. Luis explicará cómo comprobamos el cálculo.



\begin{indicacion}\emph{Relevo Benjamín \ensuremath{\rightarrow} Luis: quien sigue se acerca mientras se dice la última frase.}\end{indicacion}

\clearpage

\lamina{4}{Cómo comprobamos el cálculo}{Luis}{2:38--3:31}

\begin{indicacion}\emph{Al entrar en pantalla:} Solo el título.\end{indicacion}

\begin{indicacion}\textbf{\ensuremath{\blacktriangleright} CLIC 1} --- aparece: IFC \ensuremath{\rightarrow} malla \ensuremath{\rightarrow} una celda con sus flechas. \emph{Cuándo:} después de «gracias, Benjamín».\end{indicacion}


Gracias, Benjamín. Transformamos el IFC en cubos de 40 centímetros. En cada celda calculamos cuánto calor almacena, cuánto intercambia con sus vecinas y con el exterior, cuánto genera la hidratación y cuánto extraen los circuitos.



\begin{indicacion}\textbf{\ensuremath{\blacktriangleright} CLIC 2} --- aparece: la ecuación del balance y los dos \ensuremath{\checkmark}. \emph{Cuándo:} antes de «esta ecuación reúne ese balance».\end{indicacion}


Esta ecuación reúne ese balance. Comprobamos la conservación de energía y la respuesta frente a un caso de difusión con solución conocida: eso verifica el método numérico.



\begin{indicacion}\textbf{\ensuremath{\blacktriangleright} CLIC 3} --- aparece: tabla de sensibilidad espacio/tiempo. \emph{Cuándo:} al decir «después afinamos la malla».\end{indicacion}


Después afinamos la malla a 20 centímetros para los 24 escenarios y, por separado, reducimos el paso máximo de 300 a 150 segundos en dos casos decisivos. El diferencial cambió hasta 0,88 grados por malla y 0,0014 por tiempo. La clasificación se mantiene; esto no calibra la refrigeración real.



\clearpage

\lamina{5}{Qué muestran las temperaturas}{Luis}{3:31--4:31}

\begin{indicacion}\emph{Al entrar en pantalla:} Título y ejes vacíos.\end{indicacion}


Veamos un escenario: fuente M4 de la campaña SS y clima de 2024.



\begin{indicacion}\textbf{\ensuremath{\blacktriangleright} CLIC 1} --- aparece: gráfico de arriba (Tmax y S25) y el rótulo del pico. \emph{Cuándo:} al terminar de presentar el escenario.\end{indicacion}


En rojo está la máxima del hormigón y en azul, la mínima de S25. La máxima llega a 46,7 grados cerca de las 128 horas y luego disminuye. Distinguimos así el calentamiento de la masa del enfriamiento del borde.



\begin{indicacion}\textbf{\ensuremath{\blacktriangleright} CLIC 2} --- aparece: gráfico de \ensuremath{\Delta}T con el límite de 20 K y la marca de 14 días. \emph{Cuándo:} cuando digas «abajo está la diferencia».\end{indicacion}


Abajo está la diferencia calculada en cada instante. Su máximo ocurre después del pico de temperatura, porque el borde se enfría mientras el interior todavía conserva calor.



\begin{indicacion}\textbf{\ensuremath{\blacktriangleright} CLIC 3} --- aparece: corte a 267 h, la resta 42,12 \ensuremath{-} 25,59 y el resultado fino. \emph{Cuándo:} antes de «el peor momento de este escenario».\end{indicacion}


El peor momento de este escenario ocurre a las 267 horas: 42,12 menos 25,59 da 16,53 grados. En el catálogo fino, el peor diferencial es 15,72. La combinación de manta hasta 504 horas y sombra posterior también cumple hasta el día 28. Matías explicará qué significa ese cumplimiento.



\begin{indicacion}\emph{Relevo Luis \ensuremath{\rightarrow} Matías: quien sigue se acerca mientras se dice la última frase.}\end{indicacion}

\clearpage

\lamina{6}{Qué significa la probabilidad}{Matías}{4:31--5:25}

\begin{indicacion}\emph{Al entrar en pantalla:} Título y ejes vacíos.\end{indicacion}

\begin{indicacion}\textbf{\ensuremath{\blacktriangleright} CLIC 1} --- aparece: las 8 curvas de calor de hidratación. \emph{Cuándo:} después de «gracias, Luis».\end{indicacion}


Gracias, Luis. Conservamos las ocho curvas completas del profesor: cada una representa una historia de generación de calor. No las promediamos ni combinamos sus valores horarios.



\begin{indicacion}\textbf{\ensuremath{\blacktriangleright} CLIC 2} --- aparece: los 3 climas y la matriz 8 \ensuremath{\times} 3. \emph{Cuándo:} cuando digas «con tres climas».\end{indicacion}


Con tres climas históricos de El Rosal, de 2023, 2024 y 2025, formamos 24 escenarios. Suponemos independencia entre material y clima, y damos a cada combinación un peso de uno sobre 24. Sin dosificaciones confirmadas, interpretamos las probetas como escenarios de modelo.



\begin{indicacion}\textbf{\ensuremath{\blacktriangleright} CLIC 3} --- aparece: las ecuaciones de falla y P\_f = 0 \%. \emph{Cuándo:} antes de definir el evento de falla.\end{indicacion}


Un escenario falla si supera 70 grados de máxima o 20 de diferencial. Si supera ambos, cuenta una sola vez. Sumamos los pesos de las fallas: da cero en el catálogo, sin certificar riesgo real cero.



\clearpage

\lamina{7}{Por qué elegimos 60 m\textsuperscript{3}/h}{Matías}{5:25--6:10}

\begin{indicacion}\emph{Al entrar en pantalla:} Solo el título.\end{indicacion}

\begin{indicacion}\textbf{\ensuremath{\blacktriangleright} CLIC 1} --- aparece: tabla de los dos ritmos y la exclusión por capacidad. \emph{Cuándo:} al empezar a hablar.\end{indicacion}


Comparamos 60 y 55 metros cúbicos por hora manteniendo la misma protección, recursos y escenarios. Ambos planes cumplen los límites hasta los 28 días.



\begin{indicacion}\textbf{\ensuremath{\blacktriangleright} CLIC 2} --- aparece: cronogramas de 11 capas y gráfico D--P\_f con la meta del 5 \%. \emph{Cuándo:} cuando digas «la diferencia es el tiempo».\end{indicacion}


La diferencia es el tiempo: a 55, la colocación dura 11,58 horas, unos 58 minutos adicionales. Esperar más no garantiza menor diferencial: también influyen el clima y la exposición.



\begin{indicacion}\textbf{\ensuremath{\blacktriangleright} CLIC 3} --- aparece: la cota D \ensuremath{\geq} V/Q y el recuadro en la fila de 60. \emph{Cuándo:} al decir «y más rápido no se puede».\end{indicacion}


Y más rápido, con estos recursos, no se puede: dividir los 636,775 metros cúbicos por la capacidad máxima adoptada de 60 da 10,61 horas. Nuestra propuesta alcanza esa cota; no estamos demostrando un óptimo para cualquier capacidad o sistema constructivo.



\clearpage

\lamina{8}{Recomendación y condiciones}{Matías}{6:10--6:45}

\begin{indicacion}\emph{Al entrar en pantalla:} Recomendación, resultado fino (47,03 \ensuremath{^\circ}C \ensuremath{\cdot} 15,72 K \ensuremath{\cdot} 0/24) y la barra de margen.\end{indicacion}


Recomendamos colocar a 60 metros cúbicos por hora y 12 grados, con nueve circuitos, última manta hasta 21 días de edad local y sombra posterior.



\begin{indicacion}\textbf{\ensuremath{\blacktriangleright} CLIC 1} --- aparece: monitoreo propuesto y la condición de recalcular. \emph{Cuándo:} cuando digas «proponemos medir».\end{indicacion}


Proponemos medir núcleo, S25, caudal y temperatura del agua. Si cambian el clima o el frío, actualizamos la predicción antes de retirar protecciones.



\begin{indicacion}\textbf{\ensuremath{\blacktriangleright} CLIC 2} --- aparece: recuadro de limitación. \emph{Cuándo:} al decir «la limitación».\end{indicacion}


La limitación principal es que refrigeración y sombra están idealizadas. El resultado respalda esta decisión dentro del modelo; falta comprobarlo en obra. Muchas gracias.



\end{document}
